\section{神经网络基础回顾}

\subsection{从逻辑回归到神经网络}

上一讲介绍了神经网络的基本概念：每个神经元本质上是一个小型的逻辑回归单元。与传统机器学习中单个逻辑回归不同的是，神经网络将这些单元\textbf{级联}起来，形成多层结构。

\begin{figure}[H]
\centering
\includegraphics[width=0.95\textwidth]{slides-images/slide-002.jpg}
\caption{多层神经网络：多个逻辑回归单元级联，允许网络自动学习中间表示\protect\footnotemark}
\end{figure}
\footnotetext{Slides 第3页。}

\begin{importantbox}{神经网络的核心优势——表示学习}
神经网络最强大的特性在于\textbf{中间层的自组织}（self-organization of intermediate representations）。网络不需要人工设计特征，而是自动学习``什么样的中间表示对最终任务有用''。这正是神经网络在大多数场景下优于传统机器学习的根本原因。
\end{importantbox}

\subsection{层的矩阵表示}

当神经网络采用规则的层状结构时，每一层的计算可以用矩阵运算表示：

\begin{figure}[H]
\centering
\includegraphics[width=0.95\textwidth]{slides-images/slide-003.jpg}
\caption{单层计算的矩阵表示：$z = Wx + b$，然后 $a = f(z)$，其中 $f$ 是逐元素应用的激活函数\protect\footnotemark}
\end{figure}
\footnotetext{Slides 第4页。}

具体来说，一层神经网络的计算包含两步：

\textbf{第一步：线性变换}
$$z = Wx + b$$

其中：
\begin{itemize}
    \item $x \in \mathbb{R}^m$：输入向量
    \item $W \in \mathbb{R}^{n \times m}$：权重矩阵
    \item $b \in \mathbb{R}^n$：偏置向量
    \item $z \in \mathbb{R}^n$：线性变换输出
\end{itemize}

\textbf{第二步：非线性激活}
$$a = f(z)$$

激活函数 $f$ 逐元素应用：$f([z_1, z_2, z_3]) = [f(z_1), f(z_2), f(z_3)]$。

\subsection{NER 示例：命名实体识别}

Manning 教授使用一个具体的 NER（命名实体识别）任务来贯穿整个讲解：判断一个上下文窗口中心的词是否为地点名称。

\begin{figure}[H]
\centering
\includegraphics[width=0.95\textwidth]{slides-images/slide-004.jpg}
\caption{NER 二分类示例：输入为 5 个词的词嵌入拼接 $x \in \mathbb{R}^{5d}$，经过隐藏层 $h = f(Wx + b)$，再经点积 $s = u^T h$ 和 sigmoid 得到概率\protect\footnotemark}
\end{figure}
\footnotetext{Slides 第5页。}

该网络的完整计算过程：
\begin{enumerate}
    \item 将 5 个词的词向量拼接为输入 $x \in \mathbb{R}^{5d}$
    \item 经过隐藏层：$h = f(Wx + b)$
    \item 点积得到分数：$s = u^T h$
    \item 通过 sigmoid 得到概率：$J_t(\theta) = \sigma(s) = \frac{1}{1 + e^{-s}}$
\end{enumerate}

\begin{knowledgebox}{为什么用词嵌入而不是 one-hot？}
One-hot 向量维度极高（词表大小）且稀疏，无法表达词与词之间的语义关系。词嵌入（word embeddings）将每个词映射为一个低维稠密向量，使得语义相近的词在向量空间中距离也近。在深度学习中，$\theta$ 不仅包含权重矩阵等传统参数，\textbf{词向量本身也是需要学习的参数}。
\end{knowledgebox}

\subsection{本章小结}

神经网络通过将多个``小型逻辑回归''级联，自动学习中间表示。每一层的计算可以简洁地表示为矩阵乘法加偏置再加非线性激活：$a = f(Wx + b)$。接下来的问题是：如何训练这个网络？答案是梯度下降，而这需要我们计算梯度。

%% ========== Part 3: 激活函数 ==========

